\DocumentMetadata{
  pdfversion=2.0,
  pdfstandard=ua-2,
  lang=en-US,
  tagging=on,
  tagging-setup={math/setup=mathml-SE,tabsorder=structure}
}
\documentclass[11pt,letterpaper]{article}
\usepackage[margin=0.68in,includefoot]{geometry}
\usepackage{newcomputermodern}
\usepackage{amsmath,amscd,mathtools}
\usepackage{tikz,adjustbox}
\providecommand{\square}{\mdlgwhtsquare}
\usepackage[hidelinks]{hyperref}
\hypersetup{
  pdftitle={Topology PhD Exam, September 1989},
  pdfauthor={University of Florida Department of Mathematics},
  pdfsubject={Graduate mathematics examination},
  pdfdisplaydoctitle=true,
  bookmarksopen=true
}
\setcounter{secnumdepth}{0}
\setlength{\parindent}{0pt}
\setlength{\parskip}{0.28em}
\setlength{\emergencystretch}{3em}
\setlength{\leftmargini}{2.15em}
\setlength{\leftmarginii}{2.65em}
\setlength{\leftmarginiii}{3em}
\pagestyle{empty}
\raggedbottom
\allowdisplaybreaks
\NewTaggingSocketPlug{block/list/label}{examproblem}{%
  \tagstructbegin{tag=Lbl}\tagstructend%
  \tagstructbegin{tag=\LBody}%
  \tagstructbegin{tag=\ProblemHeadingTag,title-o={\ProblemHeadingText}}%
  \tagmcbegin{tag=\ProblemHeadingTag,actualtext={\ProblemHeadingText}}%
  #2%
  \tagmcend%
  \tagstructend%
}
\newenvironment{examproblems}[2]{%
  \def\ProblemHeadingTag{#2}%
  \AssignTaggingSocketPlug{block/list/label}{examproblem}%
  \begin{enumerate}%
  \setcounter{enumi}{#1}%
  \renewcommand{\labelenumi}{\textbf{\theenumi.}}%
  \setlength{\topsep}{0.35em}%
  \setlength{\partopsep}{0pt}%
  \setlength{\itemsep}{0.48em}%
  \setlength{\parsep}{0.18em}%
}{\end{enumerate}}
\newenvironment{examsubparts}{%
  \AssignTaggingSocketPlug{block/list/label}{default}%
  \begin{enumerate}%
  \setlength{\topsep}{0.18em}%
  \setlength{\partopsep}{0pt}%
  \setlength{\itemsep}{0.16em}%
  \setlength{\parsep}{0.1em}%
}{\end{enumerate}}
\newenvironment{examinstructions}{%
  \par\begingroup\itshape
}{%
  \par\endgroup\vspace{0.18em}
}
\makeatletter
\renewcommand\subsection{\@startsection{subsection}{2}{0pt}%
  {-1.25ex plus -.25ex minus -.1ex}%
  {0.55ex plus .1ex}%
  {\normalfont\large\bfseries}}
\renewcommand{\@maketitle}{%
  \newpage\null\vskip -1em%
  {\footnotesize\bfseries
    \href{https://www.ufl.edu/}{University of Florida}, \href{https://math.ufl.edu/}{Department of Mathematics}\par}%
  \vskip -0.5em%
  \tagstructbegin{tag=H1,title={Topology PhD Exam, September 1989}}%
  \tagpdfparaOff%
  \tagmcbegin{tag=H1}%
    {\LARGE\bfseries\href{https://gma.math.ufl.edu/exams/topology-phd/}{Topology PhD Exam}, September 1989\par}%
  \tagmcend%
  \tagpdfparaOn%
  \tagstructend%
  \vskip 0.25em%
  \tagmcbegin{artifact}%
    \hrule height 1pt%
  \tagmcend%
  \vskip 1em%
}
\makeatother
\title{Topology PhD Exam, September 1989}
\author{}
\date{}
\begin{document}
\maketitle
\thispagestyle{empty}
\begin{examproblems}{0}{H2}
\def\ProblemHeadingText{Problem 1.}
\item
State the exact homology sequence of a pair.
\def\ProblemHeadingText{Problem 2.}
\item
Describe all coverings of the torus \(T^2=S^1\times S^1\).
\def\ProblemHeadingText{Problem 3.}
\item
State the Lefschetz fixed point theorem for compact simplicial complexes.
\def\ProblemHeadingText{Problem 4.}
\item
State Tietze’s extension theorem for normal spaces.
\def\ProblemHeadingText{Problem 5.}
\item
Let \(K=K_1\cup K_2\) be a simplicial complex with \(K_1\) and \(K_2\) subcomplexes, and let \(A=K_1\cap K_2\). State the Mayer–Vietoris sequence for this.
\def\ProblemHeadingText{Problem 6.}
\item
Let~\(X\) be the set of points in the plane which is the closure of the set \(S=\{(x,\sin(1/x))\mid 0<x\leq 1\}\). Show that there is no continuous mapping \(f:[0,1]\to X\) which is onto.
\def\ProblemHeadingText{Problem 7.}
\item
Prove that a compact Hausdorff space is normal.
\def\ProblemHeadingText{Problem 8.}
\item
Let \(f,g:X\to S^n\) be any continuous mappings. Suppose that for all \(x\in X\), \(f(x)\) and \(g(x)\) are not antipodal. Show that~\(f\) and~\(g\) are homotopic.
\def\ProblemHeadingText{Problem 9.}
\item
Use the van Kampen theorem to compute the fundamental group of the surface of genus 2.
\def\ProblemHeadingText{Problem 10.}
\item
Give an example of a space~\(X\) for which \(\pi_1(X,x_0)\) is not finitely generated.
\end{examproblems}
\end{document}
